% \chapter{Literature Review}

% \section{Introduction}
% Artificial intelligence has increasingly been applied in recruitment and selection to improve efficiency, reduce manual workload, and support data-driven decision-making. Existing work in this domain spans several important areas, including automated resume screening, AI-assisted interview evaluation, speech recognition, behavioral analysis, and fairness in algorithmic hiring. While these developments have significantly advanced the state of the art, many existing approaches focus on only one stage of the recruitment pipeline. Since NeuroHire is designed as an integrated recruitment support system, it is important to review literature from each of these related domains and examine how prior work aligns with, or differs from, the proposed system.

\section{AI in Recruitment and Selection}
Artificial intelligence has been widely explored as a means of improving recruitment workflows. Rathore discusses the impact of AI on recruitment and selection processes and explains that AI can automate repetitive tasks such as resume screening, interview scheduling, and candidate communication, thereby improving efficiency and reducing manual workload for recruiters \cite{rathore_ai_recruitment}. The study also places these developments in the broader context of electronic human resource management, where digital platforms support online recruitment and virtual selection processes.

Similarly, Sha Ri examines the application of artificial intelligence in recruitment and selection, with emphasis on resume screening, interviewing, and competency assessment \cite{sha_ri_ai_recruitment}. The study explains that natural language processing can be used to extract candidate information from resumes, while machine learning and data mining can support matching between candidates and job requirements. It further highlights the role of speech recognition and computer vision in analyzing candidate responses and non-verbal behavior during interviews.

These studies demonstrate that AI has become an important part of modern recruitment systems. However, they mainly discuss the role of AI at a general level and do not present a complete multimodal framework that combines CV analysis, interview response understanding, and behavioral evaluation in a single end-to-end system. This motivates the need for an integrated system such as NeuroHire.

\section{Bias, Fairness, and Transparency in Algorithmic Hiring}
As AI systems become more common in hiring, fairness and transparency have emerged as major concerns. Raghavan \textit{et al.} analyze publicly available claims made by vendors of algorithmic hiring tools and show that many of these systems rely on question-based assessments, video interview analysis, and game-based evaluations \cite{raghavan2020bias}. Their study identifies significant shortcomings in transparency, especially regarding model validation and the fairness-related claims made by vendors.

A major concern raised in this work is that many algorithmic hiring systems depend on historical employee data. Because such data may reflect previous organizational bias, predictive models risk reproducing unfair patterns instead of reducing them \cite{raghavan2020bias}. The authors also argue that fairness is often treated in a limited way, for example through compliance-based measures, without addressing deeper issues in the data and modeling process.

This literature is important for NeuroHire because it highlights that hiring systems should not be treated as opaque decision-makers. Instead, they should provide structured and interpretable outputs. Unlike many black-box hiring tools criticized in prior work, NeuroHire is intended as a decision-support system that combines automated analysis with human-readable score components and the possibility of human review.

\section{Speech Recognition for Interview Systems}
Speech recognition is a foundational component of AI-assisted interview analysis. Radford \textit{et al.} propose a large-scale weakly supervised speech recognition approach trained on approximately 680,000 hours of multilingual and multitask audio data \cite{radford2022whisper}. Their work demonstrates strong robustness and generalization across multiple benchmarks and shows that performance improves significantly with increasing training scale and diversity. The study also supports multilingual transcription and translation, which is especially relevant for interview settings where candidates may respond in different languages.

This work is highly relevant to NeuroHire because accurate transcription is necessary before candidate responses can be analyzed for relevance, clarity, and confidence. However, the role of speech recognition in NeuroHire goes beyond general transcription. While Radford \textit{et al.} focus on robust speech recognition itself, NeuroHire uses transcription as an enabling step for structured interview response evaluation in a recruitment context.

\section{Multimodal Interview Assessment}
A growing body of research supports the use of multimodal AI for automated interview evaluation. In \textit{Automated Analysis and Prediction of Job Interview Performance}, the authors propose a framework for quantifying verbal and non-verbal behaviors in interview settings using prosodic, lexical, and facial features \cite{naim2015jobinterview}. Their results show that higher speaking rate, fewer filler words, and more positive language are associated with better interview performance. The study also demonstrates that multimodal analysis outperforms single-modality approaches in predicting interview-related traits.

Nguyen \textit{et al.} study body communication cues in employment interviews and show that gestures, posture, and speaking status can be used to predict hirability and personality-related outcomes \cite{nguyen2013bodycommunication}. Their work emphasizes that non-verbal behavior is closely linked with speech and contributes significantly to interview outcomes. However, the study also notes that part of the analysis depends on manual annotations, which limits full automation.

Zhang \textit{et al.} further extend this area by investigating asynchronous video interviews and proposing a multimodal framework for assessing personality traits and interview performance using visual, audio, and verbal signals \cite{zhang2025avi}. Their work addresses limitations in earlier studies by using a structured interview dataset designed for realistic job application scenarios and annotated by trained evaluators. The proposed baseline combines visual, audio, and textual representations and shows that multimodal systems can improve prediction performance in interview assessment tasks.

These studies strongly support the use of multimodal signals in video interview analysis. However, most of them focus on predictive modeling, dataset creation, or benchmarking rather than developing a practical recruitment support system. NeuroHire differs by using multimodal analysis within a larger system that also includes CV analysis, job-specific screening, and structured scoring for deployment-oriented candidate evaluation.

\section{Behavioral Analysis}
Behavioral signals such as gaze direction, blinking, and head movement are often used as indicators of attention and interaction behavior. Pimplaskar \textit{et al.} propose a real-time eye blinking detection and tracking system using OpenCV and show that eye-related visual cues can be extracted effectively under real-time conditions \cite{pimplaskar2013eye}. Their work demonstrates the practical feasibility of monitoring eye activity using computer vision methods and highlights the broader relevance of eye-based cues for attention-related analysis.

Head pose is another important visual signal. Huttunen \textit{et al.} study computer vision methods for head pose estimation and show that machine learning approaches can achieve strong performance in predicting yaw and pitch from facial images despite challenges such as limited data, illumination changes, and label ambiguity \cite{huttunen_headpose}. Their work provides technical support for using head orientation as a measurable behavioral feature.

For NeuroHire, these studies justify the use of gaze-related and head-pose-related signals as part of behavioral analysis during interviews. At the same time, such cues must be interpreted carefully. In NeuroHire, they are intended to act as supportive indicators of engagement and interaction behavior rather than standalone measures of competence or honesty.

\section{Existing AI-Based Interview Systems}
Some prior systems are relatively close to the goals of NeuroHire. Pandey \textit{et al.} propose an AI-based interview system that combines speech recognition, facial expression analysis, resume parsing, and machine learning-based evaluation \cite{pandey2025aibased}. Their system uses natural language processing for content-related analysis, OpenCV and convolutional neural networks for visual analysis, and a composite scoring mechanism to generate overall interview performance scores.

This work demonstrates that combining textual, vocal, and visual analysis in recruitment systems is both feasible and useful. However, NeuroHire differs in several ways. It is designed as a more structured multimodal recruitment framework rather than a collection of loosely connected modules. In addition, NeuroHire places stronger emphasis on explainable scoring, system-level integration, and controlled use of behavioral cues within a broader recruitment support pipeline.

\section{Discussion and Relevance to NeuroHire}
The literature reviewed in this chapter shows that substantial progress has been made in AI-based recruitment, speech recognition, multimodal interview analysis, and behavioral video assessment. Prior work supports the use of AI for resume screening and recruitment automation \cite{rathore_ai_recruitment,sha_ri_ai_recruitment}, robust multilingual transcription \cite{radford2022whisper}, multimodal prediction of interview performance \cite{naim2015jobinterview,nguyen2013bodycommunication,zhang2025avi}, and extraction of visual behavioral cues such as eye movement and head pose \cite{pimplaskar2013eye,huttunen_headpose}. At the same time, literature on algorithmic hiring highlights concerns related to fairness, transparency, and interpretability \cite{raghavan2020bias}.

A key observation from the reviewed work is that many existing approaches address only one part of the hiring process at a time. Some focus on recruitment automation at a high level, some on transcription, and others on multimodal interview prediction or behavioral feature extraction. Comparatively fewer systems attempt to combine CV analysis, speech-based answer evaluation, behavioral analysis, and structured decision support in one coherent framework.

This is where NeuroHire establishes its contribution. Rather than treating candidate evaluation as a single prediction task, NeuroHire integrates multiple modalities into a unified recruitment support pipeline. It combines CV-based screening, speech transcription and response analysis, and behavioral signal extraction from interview videos, while also emphasizing explainable outputs and practical usability. In this way, NeuroHire builds upon the strengths of prior work while addressing the lack of integration among existing approaches.

